\documentclass{article}

\usepackage{../handouts}

\title{CSCI 432 Handout 14: Minimum Spanning Trees}
\date{7 April 2026}
\author{Name: \rule{3in}{0.15mm}}

\begin{document}
\maketitle

\section*{Examples}

\begin{enumerate}
    \item Draw an example of a weighted graph with two distinct minimum spanning trees.

        \textbf{Answer}
        \vspace{2in}

    \item What assumptions will remove this issue?

        \textbf{Answer}
        \vspace{2in}

    \item For the next two questions, suppose we have a weighted graph
        $G=(V,E,\omega)$ meeting the asusmptions listed above whose minimum
        spanning tree is $T$.  Let $F$ be a subgraph of $T$.

        If $e \in E \setminus F$ and both endpoints of $e$ are in the same
        component of $F$, is $e \in T$? Why or why not?

        \textbf{Answer}
        \vspace{2in}

    \item Let $E' \in E \setminus F$ be the set of edges such that their
        endpoints are in different components of $F$.  Let $e$ be the
        minimum-weight edge in $E'$. Is $e \in T$? Why or why not?

        \textbf{Answer}
        \vspace{2in}



    \item What assumptions will remove this issue?

        \textbf{Answer}
        \vspace{2in}



\end{enumerate}

\pagebreak
\section*{Bor\r{u}vka's Algorithm}

In the textbook, Dr.~Erickson describes a generic MST algorithm: We start with
$F$ as the set of vertices (a spanning forest, if you will), and add edges to
$R$ until we find a MST.

In Bor\r{u}vka's algorithm, all safe edges are added at each step.

\begin{enumerate}

    \item Walk through an example, marking all safe edges and all useless edges
        at each step. (The graph you start with should have at least eight vertices
        and 10 edges, and should satisfy the assumptions we listed on the
        first page).

        \textbf{Answer}
        \vspace{3in}

    \item What is an observation you have about the useless edges?

        \textbf{Answer}
        \vspace{3in}

    \item What is the loop invariant for this algorithm?

        \textbf{Answer}
        \vspace{3in}

\end{enumerate}


\section*{Jarnik's / Prim's Algorithm}

In Jarnik's algorithm, we have a subtree $T'$ of the MST $T'$.  We obtain this
by starting at an arbitrary vertex and adding the minimum safe edge for $T'$ at
each step.


\begin{enumerate}

    \item Walk through an example. (The graph you start with should have at least eight vertices
        and 10 edges, and should satisfy the assumptions we listed on the
        first page).

        \textbf{Answer}
        \vspace{3in}

    \item What is the loop invariant for this algorithm?

        \textbf{Answer}
        \vspace{3in}

    \item What are
        some pros/cons of using Jarnik's versus Bor\r{u}vka's algorithms?

        \textbf{Answer}
        \vspace{3in}

\end{enumerate}


\section*{Kruskal's Algorithm}

In Kruskal's algorithm, we again have a spanning forest $F$ at each stage,
initialized to each vertex is its own tree. At each step, we pick the minimum
weight edge and add it to $F$, as long as it is not useless.

\begin{enumerate}

    \item Walk through an example.
        (The graph you start with should have at least eight vertices
        and 10 edges, and should satisfy the assumptions we listed on the
        first page).

        \textbf{Answer}
        \vspace{3in}

\end{enumerate}

\end{document}
